%% Question bank chapter file - generated by bank_tools.py build (do not edit while the bank is still being built:
%% it is regenerated from the per-test data in Question Bank\data). Every question is one bankq block:
%%   \begin{bankq}[id=..., source={...}, paper=.., qno=.., type=..., subtopic={...}, answer={...}, answer_src=key|solved,
%%                 working={...}, difficulty=1-10, diff_parts={K S I C F}, diff_note={...},
%%                 group=<passage id or empty>, also_in={duplicate copies in other tests}]
%% difficulty = min(10, K+S+I+C+F-1): K concepts 1-3, S steps 1-4, I insight 0-2, C computation 0-2, F format 0-1
%%   ...question LaTeX, options, figures (\bfig / \bfigi -> img/)...
%%   \end{bankq}
%% Passages shared by several questions are bankpassage blocks; their questions carry group=<passage id>.
\documentclass[10pt]{article}
\usepackage[a4paper,margin=1.8cm]{geometry}
\usepackage[T1]{fontenc}
\usepackage{amsmath,amssymb,mathtools}
\usepackage{graphicx}
\usepackage[export]{adjustbox}
\usepackage{xcolor,needspace,keyval}
\usepackage{enumitem}
\usepackage{booktabs,array,tabularx}
\usepackage[hidelinks]{hyperref}
\setlength{\parindent}{0pt}
\setlength{\parskip}{3pt}
\makeatletter
\define@key{bq}{id}{\def\bq@id{#1}}\define@key{bq}{source}{\def\bq@src{#1}}\define@key{bq}{paper}{\def\bq@paper{#1}}
\define@key{bq}{qno}{\def\bq@qno{#1}}\define@key{bq}{type}{\def\bq@type{#1}}\define@key{bq}{subtopic}{\def\bq@sub{#1}}
\define@key{bq}{answer}{\def\bq@ans{#1}}\define@key{bq}{answer_src}{\def\bq@asrc{#1}}\define@key{bq}{working}{\def\bq@work{#1}}
\define@key{bq}{group}{\def\bq@group{#1}}\define@key{bq}{also_in}{\def\bq@also{#1}}
\define@key{bq}{difficulty}{\def\bq@diff{#1}}\define@key{bq}{diff_parts}{\def\bq@dparts{#1}}\define@key{bq}{diff_note}{\def\bq@dnote{#1}}
\newenvironment{bankq}[1][]{\def\bq@work{}\def\bq@also{}\def\bq@group{}\def\bq@diff{?}\def\bq@dparts{}\def\bq@dnote{}\setkeys{bq}{#1}%
  \par\needspace{4\baselineskip}\medskip\noindent{\small\textbf{\bq@id}\hfill\color{gray}[\bq@src, Paper \bq@paper, Q\bq@qno]
  \bq@type\ $\cdot$ \bq@sub\ $\cdot$ \textbf{D\bq@diff}}\par\nopagebreak\relax}
  {\par{\small\color{teal}Answer: \textbf{\bq@ans}\ifx\bq@work\@empty\else\ (solved: \bq@work)\fi
  \ifx\bq@also\@empty\else\par\color{gray}Also in: \bq@also\fi}\par}
\define@key{bp}{id}{\def\bp@id{#1}}\define@key{bp}{source}{\def\bp@src{#1}}\define@key{bp}{questions}{\def\bp@qs{#1}}
\newenvironment{bankpassage}[1][]{\setkeys{bp}{#1}\par\needspace{4\baselineskip}\bigskip\noindent
  {\small\textbf{Passage \bp@id}\hfill\color{gray}[\bp@src; questions \bp@qs]}\par\nopagebreak\relax\itshape}{\par}
\makeatother
\newcommand{\bfig}[2]{\par\smallskip{\centering\includegraphics[width=#1pt,max width=0.62\linewidth,max height=0.32\textheight]{img/#2}\par}\smallskip}
\newcommand{\bfigi}[2]{\includegraphics[width=#1pt,max width=\linewidth,max height=0.32\textheight,valign=c]{img/#2}}
\newcommand{\bmissing}[1]{\textit{[figure missing in the source conversion: \detokenize{#1}]}}
\newenvironment{qoptions}{\begin{itemize}[leftmargin=1.8em,itemsep=1pt,topsep=2pt,label={}]}{\end{itemize}}
\newcommand{\opt}[2]{\item[(#1)] #2}
\providecommand{\fwstart}{}\providecommand{\fwend}{}\providecommand{\qfullsplit}{}
\newcolumntype{L}[1]{>{\raggedright\arraybackslash}p{#1}}
\newcolumntype{C}[1]{>{\centering\arraybackslash}p{#1}}
\newcolumntype{R}[1]{>{\raggedleft\arraybackslash}p{#1}}
\providecommand{\sectionhdr}[1]{}\providecommand{\subsectionhdr}[1]{}


\begin{document}
{\Large\bfseries ITF --- Question Bank}\par\smallskip
{\small Mathematics chapter M14. Questions from all converted test papers, grouped by source series. Answers: printed key where it matches the paper, otherwise solved (answer\_src=solved, with working). Exact repeats across papers appear once, with also\_in listing the other copies.}\par\medskip\hrule

\section*{TOAS papers}
\begin{bankq}[id=TOAS03-P1-MATH-S1-Q01, source={TOAS - Test 03 (27 Feb 2022)}, paper=1, qno=1, type=single, subtopic={Principal values}, answer={A}, answer_src=key, difficulty=2, diff_parts={K1 S1 I1 C0 F0}, diff_note={cot^-1 x lies in (0,pi) but -e^y is negative}]
Number of possible ordered pair(s) $(x,y)$ satisfying the equation $\cot^{-1}x=-e^{y}$ is -
\begin{qoptions}
\opt{A}{0}
\opt{B}{1}
\opt{C}{2}
\opt{D}{3}
\end{qoptions}
\end{bankq}
\begin{bankq}[id=TOAS03-P1-MATH-S1-Q02, source={TOAS - Test 03 (27 Feb 2022)}, paper=1, qno=2, type=single, subtopic={Principal values}, answer={D}, answer_src=key, difficulty=3, diff_parts={K1 S2 I1 C0 F0}, diff_note={tan^-1 x + cot^-1 x = pi/2 only when tan^-1 x >= 0}]
The complete set of values of $x$ for which $|\tan^{-1}x|+|\cot^{-1}x|=\dfrac\pi2$ is valid
\begin{qoptions}
\opt{A}{$(-3,3)$}
\opt{B}{$(-1,1)$}
\opt{C}{$(-\infty,0]$}
\opt{D}{$[0,\infty)$}
\end{qoptions}
\end{bankq}
\begin{bankq}[id=TOAS03-P1-MATH-S2-Q10, source={TOAS - Test 03 (27 Feb 2022)}, paper=1, qno=10, type=multi, subtopic={Equations}, answer={A,C}, answer_src=key, difficulty=3, diff_parts={K1 S2 I0 C0 F1}, diff_note={sin^-1 argument must equal 1/2, so x^2-4x+3=0}]
If $6\sin^{-1}\left(x^2-4x+\dfrac72\right)=\pi$, then-
\begin{qoptions}
\opt{A}{$x=1$}
\opt{B}{$x=2$}
\opt{C}{$x=3$}
\opt{D}{$x=4$}
\end{qoptions}
\end{bankq}
\begin{bankq}[id=TOAS03-P1-MATH-S2-Q11, source={TOAS - Test 03 (27 Feb 2022)}, paper=1, qno=11, type=multi, subtopic={Identities}, answer={C,D}, answer_src=key, difficulty=5, diff_parts={K1 S3 I1 C0 F1}, diff_note={Simplify the composite for x>0 and x<0}]
The algebraic expression for \[f(x)=\tan\left(\sin^{-1}\left(\cos\left(\tan^{-1}\dfrac x2\right)\right)\right)\] can be-
\begin{qoptions}
\opt{A}{$\dfrac\pi2$}
\opt{B}{$\dfrac1x$}
\opt{C}{$\dfrac2x$}
\opt{D}{$-\dfrac2x$}
\end{qoptions}
\end{bankq}
\begin{bankq}[id=TOAS03-P1-MATH-S3-Q03, source={TOAS - Test 03 (27 Feb 2022)}, paper=1, qno=3, type=numerical, subtopic={Identities}, answer={6.00}, answer_src=key, difficulty=4, diff_parts={K1 S2 I1 C0 F1}, diff_note={Sum formula gives (x-7)(y-7)=50, count positive solutions}]
If $x$ and $y$ are positive integers satisfying \[\tan^{-1}\left(\dfrac1x\right)+\tan^{-1}\left(\dfrac1y\right)=\tan^{-1}\left(\dfrac17\right),\] then the number of ordered pairs of $(x,y)$ is
\end{bankq}
\begin{bankq}[id=TOAS03-P1-MATH-S3-Q04, source={TOAS - Test 03 (27 Feb 2022)}, paper=1, qno=4, type=numerical, subtopic={Identities}, answer={1.50}, answer_src=key, difficulty=5, diff_parts={K1 S3 I1 C0 F1}, diff_note={Branch values for x>=3/2 give theta=pi}]
If $\theta=\sin^{-1}\left(\dfrac{2x}{1+x^2}\right)+\cos^{-1}\left(\dfrac{1-x^2}{1+x^2}\right)$ for $x\geqslant\dfrac32$, then the absolute value of $\dfrac12\left(\dfrac{\cos\theta+\tan\theta+4}{\sec\theta}\right)$ is
\end{bankq}
\begin{bankq}[id=TOAS03-P2-MATH-S1-Q02, source={TOAS - Test 03 (27 Feb 2022)}, paper=2, qno=2, type=multi, subtopic={Equations}, answer={A,B}, answer_src=key, difficulty=3, diff_parts={K1 S2 I0 C0 F1}, diff_note={sin^-1(x^2-5x+7)=pi/2 so x^2-5x+6=0}]
The values of `$x$' satisfying the equation $\sin^{-1}(x^2-5x+7)=2\tan^{-1}1$ is (are)
\begin{qoptions}
\opt{A}{2}
\opt{B}{3}
\opt{C}{4}
\opt{D}{5}
\end{qoptions}
\end{bankq}
\begin{bankq}[id=TOAS03-P2-MATH-S1-Q03, source={TOAS - Test 03 (27 Feb 2022)}, paper=2, qno=3, type=multi, subtopic={Equations}, answer={A}, answer_src=key, difficulty=6, diff_parts={K2 S3 I1 C0 F1}, diff_note={tan^-1 equality gives 4{x}=2[x]+{x}, solve by integer part}]
Consider the equation, $\tan^{-1}(4\{x\})=\tan^{-1}(x+[x])$ where $[\ ]$ and $\{\ \}$ denotes greatest integer function and fractional part function respectively .
\begin{qoptions}
\opt{A}{Number of solutions of the equation are 2.}
\opt{B}{Sum of all real solution(s) of the equation is integer}
\opt{C}{The equation has no real solution}
\opt{D}{The product of all real solution(s) of the equation is non zero.}
\end{qoptions}
\end{bankq}
\begin{bankq}[id=TOAS03-P2-MATH-S2-Q03, source={TOAS - Test 03 (27 Feb 2022)}, paper=2, qno=3, type=numerical, subtopic={Identities}, answer={7.00}, answer_src=key, difficulty=4, diff_parts={K1 S2 I1 C0 F1}, diff_note={Minimum of t^2+(pi/2-t)^2 at t=pi/4 gives pi^2/8}]
If the least value of $f(x)=(\tan^{-1}x)^2+(\cot^{-1}x)^2$ where $a$ and $b$ are coprime is $\dfrac{a\pi^2}b$, then find the value of $|a-b|$ ?
\end{bankq}
\begin{bankq}[id=TOAS03-P2-MATH-S2-Q04, source={TOAS - Test 03 (27 Feb 2022)}, paper=2, qno=4, type=numerical, subtopic={Identities}, answer={2.50}, answer_src=key, difficulty=8, diff_parts={K2 S3 I2 C1 F1}, diff_note={Reduce each inverse trig to the angle 2 tan^-1(2x/3) with branch check for x=-17/13}]
If $2\sin^{-1}\left(\dfrac{9-4x^2}{9+4x^2}\right)+3\cos^{-1}\left(\dfrac{12x}{9+4x^2}\right)+\tan^{-1}\left(\dfrac{12x}{9-4x^2}\right)=\lambda\dfrac\pi2$ for $x=\dfrac{-17}{13}$, then $\dfrac\lambda2$ is
\end{bankq}
\begin{bankq}[id=TOAS03-P2-MATH-S3-Q03, source={TOAS - Test 03 (27 Feb 2022)}, paper=2, qno=3, type=numerical, subtopic={Equations}, answer={2}, answer_src=key, difficulty=7, diff_parts={K2 S4 I1 C0 F1}, diff_note={Case x>=0 gives 6x^2-5x+1=0, x<0 impossible; 2 alpha+3 beta=2}]
If $\alpha,\beta$ ($\alpha>\beta$) are two solutions of equation $\tan^{-1}x+\cot^{-1}(-|x|)=2\tan^{-1}6x$, then $2\alpha+3\beta$ is equal to
\end{bankq}
\begin{bankq}[id=TOAS03-P2-MATH-S3-Q04, source={TOAS - Test 03 (27 Feb 2022)}, paper=2, qno=4, type=numerical, subtopic={Principal values}, answer={2}, answer_src=key, difficulty=5, diff_parts={K2 S2 I1 C0 F1}, diff_note={Range [0,pi/2) needs quadratic minimum 0: discriminant zero, sum of k}]
If range of the function $f(x)=\tan^{-1}(3x^2+(k-1)x+3)$ is $\left[0,\dfrac\pi2\right)$ for $x\in\mathbb R$, then find the sum of all possible real value(s) of $k$ ?
\end{bankq}
\begin{bankq}[id=TOAS04-P1-MATH-S2-Q06, source={TOAS - Test 04 (15 Mar 2022)}, paper=1, qno=6, type=multi, subtopic={Equations}, answer={A,B}, answer_src=key, difficulty=7, diff_parts={K2 S3 I2 C0 F1}, diff_note={The arcsine terms add to pi/2; sqrt(2-x) must lie in [0,pi/2): x>2-pi^2/4}]
Which of the following satisfy the equation $\sin^{-1}\sqrt{\dfrac{1+x}2}-\sqrt{2-x}=\cot^{-1}\left(\tan\sqrt{2-x}\right)-\sin^{-1}\sqrt{\dfrac{1-x}2}$
\begin{qoptions}
\opt{A}{$\left(2-\dfrac{\pi^2}4,1\right]$}
\opt{B}{$\left(2-\dfrac{\pi^2}4,0\right)$}
\opt{C}{$[-1,1]$}
\opt{D}{$\left[1-\dfrac{\pi^2}4,1\right]$}
\end{qoptions}
\end{bankq}
\begin{bankq}[id=TOAS04-P1-MATH-S2-Q10, source={TOAS - Test 04 (15 Mar 2022)}, paper=1, qno=10, type=multi, subtopic={Identities}, answer={B,D}, answer_src=key, difficulty=8, diff_parts={K2 S3 I2 C1 F1}, diff_note={cot^-1(n^2+3/4) telescopes as tan^-1(n+1/2)-tan^-1(n-1/2): sum is tan^-1 2}]
The sum of infinite terms of series \[\begin{aligned}{}&\cot^{-1}\left(1^2+\sin^2\dfrac{2\pi}3\right)+\cot^{-1}\left(2^2+\sin^2\dfrac{4\pi}3\right)\\&\quad+\cot^{-1}\left(3^2+\sin^2\dfrac{8\pi}3\right)+\ldots\end{aligned}\] is
\begin{qoptions}
\opt{A}{$\cot^{-1}(1)$}
\opt{B}{$\tan^{-1}(2)$}
\opt{C}{$\dfrac12\sin^{-1}\left(\dfrac45\right)$}
\opt{D}{$\cot^{-1}(1)+\cot^{-1}(3)$}
\end{qoptions}
\end{bankq}
\begin{bankpassage}[id=TOAS04-P1-MATH-S3-G3-4, source={TOAS - Test 04 (15 Mar 2022)}, questions={3-4}]
Let $\alpha$ be a real root of $f(x)=8x^3+12x^2+6x+1$, then
\end{bankpassage}
\begin{bankq}[id=TOAS04-P1-MATH-S3-Q03, source={TOAS - Test 04 (15 Mar 2022)}, paper=1, qno=3, type=numerical, subtopic={Principal values}, answer={0.00}, answer_src=key, difficulty=5, diff_parts={K2 S2 I1 C0 F1}, diff_note={8x^3+12x^2+6x+1=(2x+1)^3 so alpha=-1/2, [alpha]=-1}, group=TOAS04-P1-MATH-S3-G3-4]
The value of $\sin^{-1}(\sin[\alpha])+\cos^{-1}(\cos[\alpha])$ is equal to
\end{bankq}
\begin{bankq}[id=TOAS04-P1-MATH-S3-Q04, source={TOAS - Test 04 (15 Mar 2022)}, paper=1, qno=4, type=numerical, subtopic={Principal values}, answer={1.00}, answer_src=key, difficulty=4, diff_parts={K1 S2 I1 C0 F1}, diff_note={f is a perfect cube so one distinct real root}, group=TOAS04-P1-MATH-S3-G3-4]
The number of distinct real roots of equation $f(x)=0$, is ([.] denotes the greatest integer function)
\end{bankq}
\begin{bankq}[id=TOAS04-P1-MATH-S4-Q08, source={TOAS - Test 04 (15 Mar 2022)}, paper=1, qno=8, type=numerical, subtopic={Equations}, answer={4}, answer_src=key, difficulty=7, diff_parts={K2 S3 I2 C0 F1}, diff_note={RHS equals (x-1)^2+(y-1)^2+pi so equality forces x=y=1}]
If the solution set of $\sin^{-1}x+\sin^{-1}y\geq(x+y)^2-2[x(y+1)-(1-y)]+\pi$ is $(x_0,y_0)$ then $3x_0+y_0$ is
\end{bankq}
\begin{bankq}[id=TOAS04-P2-MATH-S1-Q04, source={TOAS - Test 04 (15 Mar 2022)}, paper=2, qno=4, type=multi, subtopic={Principal values}, answer={A,C,D}, answer_src=key, difficulty=6, diff_parts={K2 S3 I1 C0 F1}, diff_note={Graph of min/max of arcsine and arccosine; count solutions of three equations}]
Let \[f(x)=\left\{\begin{array}{ll}\min\left\{\sin^{-1}x,\cos^{-1}x\right\}&;\ 0\le x\le1\\ \max.\left\{\sin^{-1}x,\cos^{-1}x\right\}&;\ -1\le x<0\end{array}\right.\] and $f(x)=\dfrac12$, $f(x)=\dfrac14$ \& $f(x)=2$ has number of solution(s) $\lambda_1,\lambda_2,\lambda_3$ respectively, then
\begin{qoptions}
\opt{A}{$\lambda_1=2$}
\opt{B}{$\lambda_2=1$}
\opt{C}{$\lambda_2=2$}
\opt{D}{$\lambda_3=1$}
\end{qoptions}
\end{bankq}
\begin{bankq}[id=TOAS04-P2-MATH-S1-Q05, source={TOAS - Test 04 (15 Mar 2022)}, paper=2, qno=5, type=multi, subtopic={Principal values}, answer={A,B,C,D}, answer_src=key, difficulty=6, diff_parts={K2 S3 I1 C0 F1}, diff_note={f(x)=cos^-1(-{-x}) is periodic and many-one, lies in (pi/2, pi]}]
Let $f:R\to R$ defined by $f(x)=\cos^{-1}(-\{-x\})$, where $\{x\}$ is fractional part function, then which of the following is/are true ?
\begin{qoptions}
\opt{A}{$f(x)$ is many one function}
\opt{B}{period of $f(x)$ is 1}
\opt{C}{graph of $f(x)$ lies above $x$-axis}
\opt{D}{$f(x)$ is neither even nor odd function}
\end{qoptions}
\end{bankq}
\begin{bankpassage}[id=TOAS04-P2-MATH-S3-G1-2, source={TOAS - Test 04 (15 Mar 2022)}, questions={1-2}]
Consider $f(x)=\tan^{-1}\left(\dfrac{|x|}{\sqrt{1-x^2}}\right)+\cos^{-1}\left(\dfrac1{\sqrt{1+x^2}}\right)$ then
\end{bankpassage}
\begin{bankq}[id=TOAS04-P2-MATH-S3-Q01, source={TOAS - Test 04 (15 Mar 2022)}, paper=2, qno=1, type=numerical, subtopic={Principal values}, answer={3.00}, answer_src=key, difficulty=6, diff_parts={K2 S3 I1 C0 F1}, diff_note={f = sin^-1|x| + tan^-1|x| on |x|<1: range [0, 3 pi/4)}, group=TOAS04-P2-MATH-S3-G1-2]
Number of integers in the range of $f(x)$ is
\end{bankq}
\begin{bankq}[id=TOAS04-P2-MATH-S3-Q02, source={TOAS - Test 04 (15 Mar 2022)}, paper=2, qno=2, type=numerical, subtopic={Principal values}, answer={0.00}, answer_src=key, difficulty=4, diff_parts={K1 S2 I1 C0 F1}, diff_note={f increasing in |x| so f(sin t) > f(cos t) on (pi/4, pi/2): no solutions}, group=TOAS04-P2-MATH-S3-G1-2]
For $\theta\in\left(\dfrac\pi4,\dfrac\pi2\right)$, number of solution satisfying $f(\sin\theta)<f(\cos\theta)$ is
\end{bankq}
\begin{bankq}[id=TOAS05-P1-MATH-S1-Q06, source={TOAS - Test 05 (03 Apr 2022)}, paper=1, qno=6, type=multi, subtopic={Equations}, answer={B,C}, answer_src=key, difficulty=8, diff_parts={K2 S4 I1 C1 F1}, diff_note={Branches of the three inverse trigonometric functions for different intervals of x}]
For the equation $\cos^{-1}\left(\dfrac{x^2-1}{x^2+1}\right)+\sin^{-1}\left(\dfrac{2x}{x^2+1}\right)+\tan^{-1}\left(\dfrac{2x}{x^2-1}\right)=\dfrac{2\pi}3$, which of the following is/are correct.
\begin{qoptions}
\opt{A}{The number of solution is 2}
\opt{B}{It has at least one irrational solution}
\opt{C}{The sum of solution is positive}
\opt{D}{The sum of solution is negative}
\end{qoptions}
\end{bankq}
\begin{bankq}[id=TOAS05-P1-MATH-S3-Q01, source={TOAS - Test 05 (03 Apr 2022)}, paper=1, qno=1, type=numerical, subtopic={Identities}, answer={0.14}, answer_src=key, difficulty=7, diff_parts={K1 S3 I2 C1 F1}, diff_note={Terms equal tan^-1(k+1+1/7)-tan^-1(k+1/7): telescopes to cot^-1(1/7)}]
Let $S_n=\displaystyle\sum_{k=0}^{n}\tan^{-1}\dfrac{49}{49k^2+63k+57}$ then $\cot\left(\displaystyle\lim_{n\to\infty}S_n\right)$ is equal to
\end{bankq}
\begin{bankq}[id=TOAS05-P2-MATH-S1-Q02, source={TOAS - Test 05 (03 Apr 2022)}, paper=2, qno=2, type=single, subtopic={Principal values}, answer={A}, answer_src=key, difficulty=6, diff_parts={K2 S3 I2 C0 F0}, diff_note={Compare powers with bases and exponents in (0,1) for x < 1/a^a}]
Let $t_1=\left(\tan^{-1}x\right)^{\tan^{-1}x}$, $t_2=\left(\tan^{-1}x\right)^{\cot^{-1}x}$, $t_3=\left(\cot^{-1}x\right)^{\tan^{-1}x}$, $t_4=\left(\cot^{-1}x\right)^{\cot^{-1}x}$. If $x\in\left(0,\dfrac1{a^a}\right)$ where $a=2015$, then which of the following is(are) correct-
\begin{qoptions}
\opt{A}{$t_4>t_3>t_1>t_2$}
\opt{B}{$t_4>t_3$ and $t_1=t_2$}
\opt{C}{$t_1=t_2=t_3=t_4$}
\opt{D}{$t_2>t_1$ and $t_3=t_4$}
\end{qoptions}
\end{bankq}
\begin{bankq}[id=TOAS05-P2-MATH-S1-Q05, source={TOAS - Test 05 (03 Apr 2022)}, paper=2, qno=5, type=single, subtopic={Identities}, answer={B}, answer_src=key, difficulty=7, diff_parts={K2 S3 I2 C1 F0}, diff_note={Sum of arcsines equal to pi gives angle relations among x, y, z}]
If $\sin^{-1}x+\sin^{-1}y+\sin^{-1}z=\pi$, then-
\begin{qoptions}
\opt{A}{$\left(1+\sqrt{1-x^2}\right)\left(1+\sqrt{1-y^2}\right)\left(1+\sqrt{1-z^2}\right)=1$}
\opt{B}{$x\left(\sqrt{1-x^2}\right)+y\left(\sqrt{1-y^2}\right)+z\left(\sqrt{1-z^2}\right)=2xyz$}
\opt{C}{$x^2+y^2+z^2=\left(\sqrt{\left(1-x^2\right)\left(1-y^2\right)\left(1-z^2\right)}\right)$}
\opt{D}{$x+y+z=xyz$}
\end{qoptions}
\end{bankq}
\begin{bankq}[id=TOAS05-P2-MATH-S1-Q06, source={TOAS - Test 05 (03 Apr 2022)}, paper=2, qno=6, type=single, subtopic={Equations}, answer={A}, answer_src=key, difficulty=6, diff_parts={K2 S3 I1 C1 F0}, diff_note={Right-hand side equals pi: tan^-1 sum with cot^-1 gives x}]
If $\tan^{-1}(x+1)+\cot^{-1}(x-1)=\sin^{-1}\dfrac45+\cos^{-1}\dfrac35$, then $x$ has the value-
\begin{qoptions}
\opt{A}{$4\sqrt{\dfrac37}$}
\opt{B}{$4\sqrt{\dfrac73}$}
\opt{C}{$14\sqrt3$}
\opt{D}{$6\sqrt7$}
\end{qoptions}
\end{bankq}
\begin{bankq}[id=TOAS05-P2-MATH-S1-Q07, source={TOAS - Test 05 (03 Apr 2022)}, paper=2, qno=7, type=single, subtopic={Equations}, answer={A}, answer_src=key, difficulty=4, diff_parts={K1 S3 I1 C0 F0}, diff_note={(tan^-1 x)^2+(pi/2-tan^-1 x)^2=5 pi^2/8 solved for tan^-1 x}]
Number of solution(s) of $\left(\tan^{-1}x\right)^2+\left(\cot^{-1}x\right)^2=\dfrac{5\pi^2}8$, is-
\begin{qoptions}
\opt{A}{1}
\opt{B}{2}
\opt{C}{3}
\opt{D}{4}
\end{qoptions}
\end{bankq}
\begin{bankq}[id=TOAS05-P2-MATH-S2-Q01, source={TOAS - Test 05 (03 Apr 2022)}, paper=2, qno=1, type=numerical, subtopic={Identities}, answer={7}, answer_src=key, difficulty=6, diff_parts={K2 S3 I1 C0 F1}, diff_note={Both inverse trigonometric parts reduce to complementary angles: f(x)=1}]
Let {\small\[f(x)=\cot\left(\dfrac{\cos^{-1}\sin\left(\cos^{-1}x\right)+\sin^{-1}\left(\cos\left(\sin^{-1}x\right)\right)}{2}\right),\]} then find the value of $\displaystyle\sum_{x=1}^{7}f\left(\dfrac x8\right)$.
\end{bankq}
\begin{bankq}[id=TOAS05-P2-MATH-S2-Q02, source={TOAS - Test 05 (03 Apr 2022)}, paper=2, qno=2, type=numerical, subtopic={Principal values}, answer={2}, answer_src=key, difficulty=6, diff_parts={K2 S3 I1 C0 F1}, diff_note={Range of sin cot^-1 cos tan^-1 x}]
The equation $\sin\cot^{-1}\cos\tan^{-1}x=a$ has real solution iff complete range of $a$ is $[p,q]$. Value of $\left[q+\dfrac1p\right]$ is (where $[.]$ denotes greatest integer function)
\end{bankq}
\begin{bankq}[id=TOAS05-P2-MATH-S2-Q03, source={TOAS - Test 05 (03 Apr 2022)}, paper=2, qno=3, type=numerical, subtopic={Identities}, answer={1}, answer_src=solved, working={$\sin^{-1}\frac1{\sqrt{10}}=\tan^{-1}\frac13$, $\tan^{-1}\frac29$, $\tan^{-1}\frac4{33}$: general term $\tan^{-1}2^n-\tan^{-1}2^{n-1}$, sum $=\frac\pi2-\frac\pi4=\frac\pi4$, $\tan\alpha=1$}, difficulty=7, diff_parts={K1 S3 I2 C1 F1}, diff_note={Terms are tan^-1(1/3), tan^-1(2/9), ...: general term tan^-1 2^n - tan^-1 2^{n-1}}]
If $\sin^{-1}\dfrac1{\sqrt{10}}+\sin^{-1}\dfrac2{\sqrt{85}}+\sin^{-1}\left(\dfrac4{\sqrt{1105}}\right)+\ldots\ldots\infty=\alpha$, then value of $\tan\alpha$ is
\end{bankq}
\begin{bankq}[id=TOAS05-P2-MATH-S2-Q04, source={TOAS - Test 05 (03 Apr 2022)}, paper=2, qno=4, type=numerical, subtopic={Identities}, answer={3}, answer_src=solved, working={$\tan^{-1}x+\tan^{-1}y+\tan^{-1}z=\frac\pi2$ gives $xy+yz+zx=1$; with $x+y+z=\sqrt3$ the squares sum to 1 so $x=y=z=\frac1{\sqrt3}$ and $\frac1{xyz}=3\sqrt3=k\sqrt k$, $k=3$}, difficulty=7, diff_parts={K2 S3 I2 C0 F1}, diff_note={tan^-1 sum is pi/2 so xy+yz+zx=1; with x+y+z=sqrt3 all equal 1/sqrt3}]
If $\tan^{-1}x+\tan^{-1}y=\cot^{-1}z$ is such that $x+y+z=\sqrt3$, then value of $\dfrac1{xyz}$ can be expressed as $k\sqrt k$, where $k$ is
\end{bankq}
\begin{bankq}[id=TOAS06-P1-MATH-S2-Q05, source={TOAS - Test 06 (17 Apr 2022)}, paper=1, qno=5, type=multi, subtopic={Equations}, answer={A,C,D}, answer_src=key, difficulty=8, diff_parts={K2 S4 I1 C1 F1}, diff_note={Count roots of the three inverse-trigonometric equations graphically}]
If the number of solutions of equations $\sin^{-1}x+|x|=1$, $\tan^{-1}x+|x|=1$ and $\sec^{-1}x+|x|=1$ are $n_1$, $n_2$ and $n_3$ respectively, then
\begin{qoptions}
\opt{A}{$n_1=1$}
\opt{B}{$n_2=1$}
\opt{C}{$n_2=2$}
\opt{D}{$n_3=1$}
\end{qoptions}
\end{bankq}
\begin{bankq}[id=TOAS06-P1-MATH-S3-Q03, source={TOAS - Test 06 (17 Apr 2022)}, paper=1, qno=3, type=numerical, subtopic={Identities}, answer={0.25}, answer_src=key, difficulty=7, diff_parts={K1 S3 I2 C1 F1}, diff_note={Terms equal tan^-1 2^k - tan^-1 2^{k-1}-type differences: sum and k}]
If the sum of $n$ terms of the series $\tan^{-1}3+\tan^{-1}\dfrac92+\tan^{-1}\dfrac{33}{4}+\tan^{-1}\dfrac{129}{8}+\ldots.$ is $(2n+1)k\pi-\tan^{-1}2^n$, then $k$ is
\end{bankq}
\begin{bankq}[id=TOAS07-P1-MATH-S2-Q08, source={TOAS - Test 07 (01 May 2022)}, paper=1, qno=8, type=multi, subtopic={Identities}, answer={A,B}, answer_src=key, difficulty=7, diff_parts={K1 S3 I2 C1 F1}, diff_note={Line y=mx meets the periodic triangular wave in exactly 5 points}]
For which of the following values of `$m$' the line $y=mx$ cuts the graph of $y=\sin^{-1}(\sin x)$ in exactly 5 points -
\begin{qoptions}
\opt{A}{$\dfrac15$}
\opt{B}{$-\dfrac14$}
\opt{C}{$-\dfrac18$}
\opt{D}{$-\dfrac19$}
\end{qoptions}
\end{bankq}
\begin{bankq}[id=TOAS07-P1-MATH-S2-Q10, source={TOAS - Test 07 (01 May 2022)}, paper=1, qno=10, type=multi, subtopic={Identities}, answer={A,B,D}, answer_src=key, difficulty=8, diff_parts={K2 S3 I2 C1 F1}, diff_note={Inverse trig relation gives xy+yz+zx=1 type; AM-GM inequality on x+y+z}]
If $\sec^{-1}\left(\sqrt{1+x^2}\right)+\mathrm{cosec}^{-1}\left(\dfrac{\sqrt{1+y^2}}{y}\right)+\cot^{-1}\left(\dfrac1z\right)=\pi$ $(x,y,z>0)$ then $x+y+z$ is always -
\begin{qoptions}
\opt{A}{equal to $xyz$}
\opt{B}{greater than or equal to $3(xyz)^{1/3}$}
\opt{C}{greater than or equal to $4(xyz)^{1/3}$}
\opt{D}{greater than or equal to $5(xyz)^{1/3}$}
\end{qoptions}
\end{bankq}
\begin{bankq}[id=TOAS07-P2-MATH-S1-Q02, source={TOAS - Test 07 (01 May 2022)}, paper=2, qno=2, type=multi, subtopic={Principal values}, answer={A,B,D}, answer_src=key, difficulty=7, diff_parts={K2 S3 I2 C0 F1}, diff_note={Sum of arcsines equals pi^2 forces all four equal to one: determinant values}]
If $(\sin^{-1}x+\sin^{-1}w).(\sin^{-1}y+\sin^{-1}z)=\pi^2$ then $\begin{vmatrix}x^{n_1}&y^{n_2}\\z^{n_3}&w^{n_4}\end{vmatrix}$, (where $n_1,n_2,n_3,n_4\in N$) -
\begin{qoptions}
\opt{A}{has a maximum value 2}
\opt{B}{has a minimum value $-2$}
\opt{C}{is independent of $n_i$ $(i=1,2,3,4)$}
\opt{D}{can take value 0}
\end{qoptions}
\end{bankq}
\begin{bankq}[id=TOAS07-P2-MATH-S1-Q04, source={TOAS - Test 07 (01 May 2022)}, paper=2, qno=4, type=multi, subtopic={Principal values}, answer={B,C,D}, answer_src=key, difficulty=5, diff_parts={K1 S3 I1 C0 F1}, diff_note={f(x)=2x sqrt(1-x^2) type behaviour: one-one, odd, continuity}]
Let $f:[-1,1]\to[-1,1]$, $f(x)=\sin(2\sin^{-1}x)$ \& let $g(x)=\dfrac{f(x)}{\sqrt{1-x^2}}$, then -
\begin{qoptions}
\opt{A}{$f(x)$ is one-one function}
\opt{B}{$f(x)$ is odd function}
\opt{C}{$g(x)$ is continuous in $(-1,1)$}
\opt{D}{$f(x)$ is continuous for all $x\in[-1,1]$}
\end{qoptions}
\end{bankq}
\begin{bankpassage}[id=TOAS07-P2-MATH-S3-G1-2, source={TOAS - Test 07 (01 May 2022)}, questions={1-2}]
Let $f(x)=\sin^{-1}x+\sin^{-1}(1-x)-\cos^{-1}x$, then answer the following
\end{bankpassage}
\begin{bankq}[id=TOAS07-P2-MATH-S3-Q01, source={TOAS - Test 07 (01 May 2022)}, paper=2, qno=1, type=numerical, subtopic={Principal values}, answer={2.00}, answer_src=key, difficulty=7, diff_parts={K2 S3 I1 C1 F1}, diff_note={Number of roots of sin^-1x+sin^-1(1-x)=cos^-1x on [0,1]}, group=TOAS07-P2-MATH-S3-G1-2]
Number of root(s) of the equation $f(x)=0$ is -
\end{bankq}
\begin{bankq}[id=TOAS07-P2-MATH-S3-Q02, source={TOAS - Test 07 (01 May 2022)}, paper=2, qno=2, type=numerical, subtopic={Principal values}, answer={0.00}, answer_src=key, difficulty=5, diff_parts={K1 S3 I1 C0 F1}, diff_note={G(0)=2 f(0) from the definition}, group=TOAS07-P2-MATH-S3-G1-2]
If $G(x)=f(x)+f(-x)$ then value of $G(0)$ is equal to
\end{bankq}
\begin{bankq}[id=TOAS08-P2-MATH-S3-Q03, source={TOAS - Test 08 (15 May 2022)}, paper=2, qno=3, type=numerical, subtopic={Identities}, answer={2.00}, answer_src=key, difficulty=8, diff_parts={K2 S3 I2 C1 F1}, diff_note={Inverse trigonometric simplification gives f(x)=x^2-like and a geometric sum}]
\[\begin{aligned}f(x)={}&\tan(\sin^{-1}(\cos(\sin^{-1}(x))))\\&\times\tan(\cos^{-1}(\sin(\cos^{-1}x))),\end{aligned}\]where $x\in(0,1)$ then $\displaystyle\sum_{r=0}^{\infty}\dfrac{1}{2^r}f\left(\dfrac{1}{2^r}\right)$ is equal to
\end{bankq}
\begin{bankq}[id=TOAS09-P1-MATH-S4-Q09, source={TOAS - Test 09 (30 May 2022)}, paper=1, qno=9, type=numerical, subtopic={Principal values}, answer={4}, answer_src=key, difficulty=8, diff_parts={K3 S3 I1 C1 F1}, diff_note={Principal values: standard approach for this subtopic; 3 concept(s), 3 step(s), needs a non-obvious insight, heavier computation}]
The value of $k$ ($k>0$) such that the length of the longest interval in which the function $f(x)=\sin^{-1}|\sin kx|+\cos^{-1}(\cos kx)$ is constant is $\dfrac{\pi}{4}$ will be
\end{bankq}
\begin{bankq}[id=TOAS09-P2-MATH-S1-Q01, source={TOAS - Test 09 (30 May 2022)}, paper=2, qno=1, type=single, subtopic={Identities}, answer={D}, answer_src=key, difficulty=2, diff_parts={K1 S2 I0 C0 F0}, diff_note={Identities: standard approach for this subtopic; 1 concept(s), 2 step(s)}]
$\displaystyle\lim_{n\to\infty}\dfrac{\displaystyle\sum_{r=0}^{n}\tan^{-1}\left(1+r+r^2\right)}{n}$ is equal to-
\begin{qoptions}
\opt{A}{1}
\opt{B}{2}
\opt{C}{$\dfrac\pi4$}
\opt{D}{$\dfrac\pi2$}
\end{qoptions}
\end{bankq}
\begin{bankq}[id=TOAS09-P2-MATH-S1-Q17, source={TOAS - Test 09 (30 May 2022)}, paper=2, qno=17, type=single, subtopic={Principal values}, answer={D}, answer_src=key, difficulty=4, diff_parts={K2 S2 I1 C0 F0}, diff_note={Principal values: standard approach for this subtopic; 2 concept(s), 2 step(s), needs a non-obvious insight}]
If the equation $\sin^{-1}(x-1)+\cos^{-1}(x-3)+\tan^{-1}\left(\dfrac{x}{-x^2+2}\right)=m$ holds, then the value of `$m$' is
\begin{qoptions}
\opt{A}{$\dfrac{-3\pi}{4}$}
\opt{B}{$\dfrac{3\pi}{4}$}
\opt{C}{$\dfrac{-5\pi}{4}$}
\opt{D}{$\dfrac{5\pi}{4}$}
\end{qoptions}
\end{bankq}
\end{document}
